\documentclass[11pt,a4paper]{article}
\usepackage[margin=22mm]{geometry}
\usepackage{amsmath,amssymb,amsthm,booktabs,microtype}
\usepackage{xurl}
\usepackage{hyperref}
\hypersetup{
  hidelinks,
  pdftitle={Heavy-Tail First-Order Variance Correction for the RSI Log-Odds Coordinate},
  pdfauthor={Paweł Małmyga},
  pdfsubject={A technical note for symmetric innovations with a finite (3+delta) moment},
  pdfkeywords={Student-t, heavy tails, variance expansion, endpoint coordinate, RSI log odds}
}
\usepackage[T1]{fontenc}
\usepackage{lmodern}
\newtheorem{theorem}{Theorem}
\newtheorem{lemma}{Lemma}
\newtheorem{corollary}{Corollary}
\newcommand{\E}{\mathbb E}
\newcommand{\Var}{\operatorname{Var}}
\newcommand{\Prob}{\mathbb P}
\newcommand{\atanh}{\operatorname{atanh}}
\title{Heavy-Tail First-Order Variance Correction\\\large For the RSI log-odds coordinate under a finite $(3+\delta)$ moment}
\author{Pawe{\l} Ma{\l}myga\\\small Independent researcher}
\date{September 2026}
\begin{document}
\maketitle

\begin{abstract}
This note extends the first law-dependent finite-horizon variance correction
from exponentially light tails to symmetric innovation laws with a finite
$(3+\delta)$ absolute moment and a logarithmic lower-tail condition.  The
result is
\[
\Var(L_\alpha)=2m_2\alpha+\frac{15m_2^2+2m_2-8m_3}{3}\alpha^2+o(\alpha^2),
\]
and therefore
\[
g_F^*(n)=\frac{C_F^2}{\Var(L_{1/n})}=n-K_F+o(1),
\qquad
K_F=\frac{15m_2^2+2m_2-8m_3}{6m_2}.
\]
For variance-normalized Student-$t_\nu$ innovations this applies for every
$\nu>3$ and gives an explicit $K_\nu$.  The proof uses a shrinking big-jump
cutoff and termwise-truncated centered variables, so no fourth raw moment is
required when $3<p<4$.  The expansion is pointwise in the innovation law: its
constants and rates may depend on $\delta$ and $\E R^{3+\delta}$, and no
uniformity as $\delta\downarrow0$ is claimed.
\end{abstract}

\section*{Verification scope}
This is a conventional mathematical proof, not a theorem accepted by the Lean
kernel.  The accompanying Lean repository checks
the exact finite-path endpoint identities, exponential-weight algebra,
conditional log-odds transfer results, and a finite-variance central limit
theorem for normalized \emph{linear} Wilder rows.  It does not currently prove
the nonlinear variance expansion stated in this note.

The repository's pinned SymPy calculation independently reconstructs the
formal moment and cumulant coefficients and their law-specific substitutions.
It checks the finite coefficient algebra; the tail localization, indicator
removal, and analytic remainder estimates are established by the argument
below and are not machine-verified.

\section{Setup}
Let $\{\varepsilon_j\}_{j\ge0}$ be i.i.d., symmetric, nondegenerate, and satisfy $\Prob(\varepsilon=0)=0$.  Put
\[
\mu_1=\E|\varepsilon|,\qquad R=\frac{|\varepsilon|}{\mu_1},\qquad Q=\operatorname{sgn}(\varepsilon).
\]
Symmetry and absence of an atom at zero imply that $Q$ is Rademacher and independent of $R$.  Hence $\E R=1$.  For $\alpha\in(0,1)$ define geometric weights
\[
w_j=\alpha(1-\alpha)^j,\qquad \sum_{j\ge0}w_j=1,
\]
and
\[
x=\sum_{j\ge0}w_jQ_jR_j,\qquad y=\sum_{j\ge0}w_j(R_j-1).
\]
Then the signed directional balance is
\[
F_\alpha=\frac{x}{1+y},
\]
and the full RSI log-odds coordinate is
\[
L_\alpha=2\atanh(F_\alpha).
\]
The law of $L_\alpha$ is symmetric, hence $\E L_\alpha=0$ whenever $L_\alpha$ is integrable.

Write
\[
m_k=\E R^k,
\qquad
s_r(\alpha)=\sum_{j\ge0}w_j^r
=\frac{\alpha^r}{1-(1-\alpha)^r}.
\]
In particular,
\[
s_2=\frac\alpha2+\frac{\alpha^2}{4}+O(\alpha^3),
\qquad
s_3=\frac{\alpha^2}{3}+O(\alpha^3).
\]

\section{Main statement}
\begin{theorem}[First finite-horizon variance correction without exponential tails]
Assume that for some $\delta>0$
\[
\E R^{3+\delta}<\infty.
\]
Assume also that
\[
\E\bigl[(-\log R)_+^q\bigr]<\infty
\]
for some
\[
q>\frac{2(2+\delta)}{\delta}.
\]
Given such a $q$, choose $\gamma>0$ so small that
\begin{equation}\label{eq:gamma-choice}
0<\gamma<\min\left\{
\frac{\delta}{3+\delta},\,
\frac{2+\delta-2q/(q-2)}{3+\delta}
\right\}.
\end{equation}
The second upper bound is positive precisely because of the strict inequality imposed on $q$.  Then, as $\alpha\downarrow0$,
\begin{equation}\label{eq:var-exp}
\boxed{
\Var(L_\alpha)
=2m_2\alpha
+\frac{15m_2^2+2m_2-8m_3}{3}\alpha^2
+o(\alpha^2).}
\end{equation}
Consequently, with
\[
C_F^2=2m_2=\frac{2\E\varepsilon^2}{(\E|\varepsilon|)^2},
\]
the exact-variance normalizer satisfies
\begin{equation}\label{eq:g-shift}
\boxed{
 g_F^*(n):=\frac{C_F^2}{\Var(L_{1/n})}
=n-K_F+o(1),}
\end{equation}
where
\begin{equation}\label{eq:K}
\boxed{
K_F=\frac{15m_2^2+2m_2-8m_3}{6m_2}.}
\end{equation}
\end{theorem}

\paragraph{Pointwise nature of the expansion.}
For each fixed innovation law satisfying the stated assumptions, one may fix an admissible $\delta>0$, then choose $q$ and $\gamma$ as above.  The remainder $o(\alpha^2)$ is therefore pointwise in the law.  Neither the implied constants nor the rate are asserted to be uniform over a family for which $\delta\downarrow0$ or $\E R^{3+\delta}$ diverges.

The proof is split into localization, complement control, and local moment algebra.

\section{Weight and truncation estimates}
Fix
\[
p=3+\delta.
\]
Choose $\gamma$ according to \eqref{eq:gamma-choice} and define
\[
\beta=p-1-\gamma p=2+\delta-\gamma(3+\delta).
\]
Then
\[
\beta>\frac{2q}{q-2}>2,
\qquad
\beta\left(1-\frac2q\right)>2.
\]
Choose $q$ first and then a sufficiently small positive $\gamma$.  Put
\[
a_\alpha=\alpha^\gamma,
\qquad
H_\alpha=\left\{\sup_{j\ge0}w_jR_j>a_\alpha\right\}.
\]

\begin{lemma}[Master tail estimates]
For $0\le r\le p$,
\begin{equation}\label{eq:tail-master}
\sum_j w_j^r\E\!\left[R^r\mathbf 1_{\{w_jR>a_\alpha\}}\right]
\le \E R^p\,a_\alpha^{r-p}s_p
=O\!\left(\alpha^{p-1-\gamma(p-r)}\right).
\end{equation}
For $r\ge p$,
\begin{equation}\label{eq:trunc-master}
\sum_j w_j^r\E\!\left[R^r\mathbf 1_{\{w_jR\le a_\alpha\}}\right]
\le \E R^p\,a_\alpha^{r-p}s_p
=O\!\left(\alpha^{p-1+\gamma(r-p)}\right).
\end{equation}
In particular,
\begin{equation}\label{eq:H-prob}
\Prob(H_\alpha)=O(\alpha^\beta)=o(\alpha^2).
\end{equation}
Moreover, the lost first, second, and third weighted moments caused by truncation are all $o(\alpha^2)$.
\end{lemma}

\begin{proof}
On $\{R>a/w_j\}$ and for $r\le p$,
\[
R^r\le R^p(a/w_j)^{r-p}.
\]
For $r\ge p$ the same inequality holds on $\{R\le a/w_j\}$ with the inequality direction appropriate to $r-p\ge0$.  Summation gives \eqref{eq:tail-master}--\eqref{eq:trunc-master}.  Setting $r=0$ and using a union bound gives \eqref{eq:H-prob}.  For $r=1,2,3$, the three exponents in \eqref{eq:tail-master} are, respectively,
\[
2+\delta-\gamma(2+\delta),\qquad
2+\delta-\gamma(1+\delta),\qquad
2+\delta-\gamma\delta.
\]
Each is strictly larger than $2$ under
$\gamma<\delta/(3+\delta)$.  Hence the corresponding lost weighted moments
are $o(\alpha^2)$.
\end{proof}

\section{Localization without an exponential moment}
Define termwise-truncated magnitudes
\[
\widetilde R_j=R_j\mathbf 1_{\{w_jR_j\le a_\alpha\}},
\]
and centered truncated variables
\[
A_j=Q_j\widetilde R_j,
\qquad
B_j=\widetilde R_j-\E\widetilde R_j.
\]
Set
\[
\widetilde x=\sum_jw_jA_j,
\qquad
z=\sum_jw_jB_j,
\qquad
b_\alpha=\sum_jw_j(\E\widetilde R_j-1).
\]
By symmetry $\E A_j=0$, while Lemma 1 with $r=1$ gives
\begin{equation}\label{eq:bias}
|b_\alpha|
\le \sum_jw_j\E[R_j\mathbf 1_{\{w_jR_j>a_\alpha\}}]
=o(\alpha^2).
\end{equation}
On $H_\alpha^c$ we have exactly
\begin{equation}\label{eq:trunc-identities}
x=\widetilde x,
\qquad
y=z+b_\alpha.
\end{equation}
The event $H_\alpha^c$ bounds each weighted summand $w_jR_j$ by $a_\alpha$;
it does not imply $|x|\le a_\alpha$ or $|y|\le a_\alpha$.  The infinite sums
are controlled below through variance and moment bounds.

\begin{lemma}[Uniform moments of the truncated vector]\label{lem:trunc-moments}
For every fixed $r\ge2$,
\[
\sup_{0<\alpha<\alpha_0}
\left(\E|\widetilde x|^r+\E|z|^r\right)<\infty.
\]
More precisely, Rosenthal's inequality bounds each moment by a constant multiple of a variance term of order $O(\alpha^{r/2})$ plus a diagonal term.  For $r\le p$ the diagonal term is $O(s_r)=O(\alpha^{r-1})$; for $r>p$ Lemma 1 gives
\[
\sum_jw_j^r\E\widetilde R_j^r
=O\!\left(\alpha^{p-1+\gamma(r-p)}\right).
\]
The same conclusion holds for $B_j$ after centering.
\end{lemma}

\begin{proof}
Apply Rosenthal's inequality first to finite partial sums.  The sum of variances is $O(\alpha)$ because $\sum_jw_j^2=O(\alpha)$ and the truncated second moments are uniformly bounded.  The stated diagonal bounds follow from the finite $p$-moment when $r\le p$ and from \eqref{eq:trunc-master} when $r>p$.  Letting the number of summands tend to infinity gives the result by $L^r$ convergence.
\end{proof}

The centered weighted summands $w_jA_j$ have absolute size at most $a_\alpha$, and $w_jB_j$ have absolute size at most $a_\alpha+O(\alpha)=O(a_\alpha)$.  Their total variances are $O(\alpha)$.

\begin{lemma}[Shrinking-cutoff concentration]
For each fixed sufficiently small $\eta>0$ there are constants $c,C>0$ such that
\begin{equation}\label{eq:local-prob}
\Prob\bigl((|x|>\eta\ \text{or}\ |y|>\eta)\cap H_\alpha^c\bigr)
\le C\exp(-c\alpha^{-\gamma})
\end{equation}
for all sufficiently small $\alpha$.  Consequently, if
\[
G_\alpha=H_\alpha^c\cap\{|x|\le\eta,\ |y|\le\eta\},
\]
then
\begin{equation}\label{eq:Gcomp}
\Prob(G_\alpha^c)=O(\alpha^\beta)+O(e^{-c\alpha^{-\gamma}}).
\end{equation}
\end{lemma}

\begin{proof}
Apply Bernstein's inequality to finite partial sums of the independent centered variables $w_jA_j$ and $w_jB_j$.  The variance proxy is $O(\alpha)$ and the maximal summand is $O(a_\alpha)=O(\alpha^\gamma)$.  For a fixed deviation $u>0$,
\[
2\exp\left[-c\min\left(\frac{u^2}{\alpha},\frac{u}{\alpha^\gamma}\right)\right]
\le 2e^{-c'\alpha^{-\gamma}},
\]
since $0<\gamma<1$.  Pass to the infinite sums by $L^2$ convergence of the partial sums.  By \eqref{eq:bias}, $y=z+b_\alpha$ on $H_\alpha^c$ with $b_\alpha=o(1)$, so the deterministic bias is absorbed into $u$ for small $\alpha$.  Add \eqref{eq:H-prob}.
\end{proof}

\paragraph{Why the cutoff must shrink.}
A fixed cutoff $a_\alpha\equiv a>0$ would only make the Bernstein bounded-summand term $\exp(-c/a)$, which does not vanish with $\alpha$.  The shrinking cutoff $a_\alpha=\alpha^\gamma$ is therefore essential in the polynomial-tail proof.

\section{The complement contributes below order
  \texorpdfstring{$\alpha^2$}{alpha squared}}
Let $T_+$ and $T_-$ be the first indices with $Q_j=+1$ and $Q_j=-1$, respectively.  They are geometric with success probability $1/2$ and are independent of the magnitudes.  Since
\[
U\ge \mu_1w_{T_+}R_{T_+},\qquad D\ge \mu_1w_{T_-}R_{T_-},
\]
we obtain
\[
(-\log U)_+
\le C+\log(1/\alpha)+T_+[-\log(1-\alpha)]+(-\log R_{T_+})_+,
\]
and similarly for $D$.  Positive logarithms are uniformly controlled by the finite $p$-moment, because $U,D\le \mu_1\sum_jw_jR_j$ and Jensen's inequality gives a uniform positive moment of the weighted average.

\begin{lemma}[Logarithmic moments and bad-set negligibility]
Under the assumptions of Theorem 1,
\begin{equation}\label{eq:Lq}
\E|L_\alpha|^q=O\bigl(1+\log^q(1/\alpha)\bigr).
\end{equation}
Moreover,
\begin{equation}\label{eq:badL2}
\E\left[L_\alpha^2\mathbf 1_{G_\alpha^c}\right]=o(\alpha^2).
\end{equation}
\end{lemma}

\begin{proof}
The first-success bounds above, the finite $q$-moment of $(-\log R)_+$, and the geometric moments of $T_\pm$ yield \eqref{eq:Lq}.  By H\"older and \eqref{eq:Gcomp},
\[
\E[L_\alpha^2\mathbf 1_{G_\alpha^c}]
\le (\E|L_\alpha|^q)^{2/q}\Prob(G_\alpha^c)^{1-2/q}.
\]
The polynomial contribution is
$O(\log^2(1/\alpha)\alpha^{\beta(1-2/q)})=o(\alpha^2)$ by construction.
The stretched-exponential contribution is smaller than every power of
$\alpha$.
\end{proof}

\section{Local Taylor algebra and moment bookkeeping}
Choose $\eta<1/3$.  On $G_\alpha$,
\[
\left|\frac{x}{1+y}\right|\le \frac{\eta}{1-\eta}<\frac12,
\]
so the Taylor series is uniformly regular.  Direct expansion gives
\begin{equation}\label{eq:taylor}
\frac{L_\alpha^2}{4}
=x^2-2x^2y+3x^2y^2+\frac23x^4+\mathcal R_5,
\end{equation}
where on $G_\alpha$
\[
|\mathcal R_5|\le C_\eta(|x|+|y|)^5.
\]

\paragraph{Why total degree four is sufficient.}
On the fixed neighborhood defining $G_\alpha$, the map
$2\atanh(x/(1+y))$ is analytic and the displayed polynomial contains every
term in its square through total degree four.  Terms of total degree at least
five are therefore absorbed into $\mathcal R_5$.  Consequently, no such term
can affect the coefficient of $\alpha^2$ once
\[
\E\bigl[(|x|+|y|)^5\mathbf 1_{G_\alpha}\bigr]=o(\alpha^2)
\]
has been established.  Lemma 5 proves precisely this analytic remainder
condition; the symbolic calculation checks only the resulting finite polynomial
algebra.

\begin{lemma}[Local moments through total degree four]
Under the assumptions of Theorem 1,
\begin{align}
\E[x^2\mathbf 1_{G_\alpha}]&=m_2s_2+o(\alpha^2),\label{eq:m20}\\
\E[x^2y\mathbf 1_{G_\alpha}]&=(m_3-m_2)s_3+o(\alpha^2),\label{eq:m21}\\
\E[x^4\mathbf 1_{G_\alpha}]&=3m_2^2s_2^2+o(\alpha^2),\label{eq:m40}\\
\E[x^2y^2\mathbf 1_{G_\alpha}]&=m_2(m_2-1)s_2^2+o(\alpha^2),\label{eq:m22}\\
\E[(|x|+|y|)^5\mathbf 1_{G_\alpha}]&=o(\alpha^2).\label{eq:m5}
\end{align}
\end{lemma}

\begin{proof}
We work with the centered termwise-truncated vector $(\widetilde x,z)$ from
Section 4.  This does not require an untruncated fourth moment when $3<p<4$.

\paragraph{Removing the indicators.}
Let $P$ be any fixed polynomial in two variables of total degree at most five.  Lemma~\ref{lem:trunc-moments} gives uniformly bounded moments of $P(\widetilde x,z)$ of arbitrarily high fixed order.  Since $\beta>2$, choose $s>1$ so large that $\beta(1-1/s)>2$.  H\"older therefore gives
\[
\E\bigl[|P(\widetilde x,z)|\mathbf 1_{H_\alpha}\bigr]
\le
\bigl(\E|P(\widetilde x,z)|^s\bigr)^{1/s}
\Prob(H_\alpha)^{1-1/s}
=o(\alpha^2).
\]
On $H_\alpha^c$, identities \eqref{eq:trunc-identities} hold.  The additional localization failure
$\{|x|>\eta\text{ or }|y|>\eta\}\cap H_\alpha^c$ has stretched-exponentially small probability by Lemma 3, and another H\"older bound using Lemma~\ref{lem:trunc-moments} makes its contribution $o(\alpha^M)$ for every fixed $M$.  Thus, for the polynomials below, expectations localized to $G_\alpha$ may be replaced, up to $o(\alpha^2)$, by full expectations of the corresponding expressions in $(\widetilde x,z+b_\alpha)$.

\paragraph{Second and third degree.}
Put
\[
V_A=\sum_jw_j^2\E\widetilde R_j^2,
\qquad
V_B=\sum_jw_j^2\Var(\widetilde R_j).
\]
The $r=1,2$ tail bounds in Lemma 1 imply
\begin{equation}\label{eq:VA-VB}
V_A=m_2s_2+o(\alpha^2),
\qquad
V_B=(m_2-1)s_2+o(\alpha^2).
\end{equation}
Hence $\E\widetilde x^2=V_A$.  Because the signs are independent and centered, every term in $\E(\widetilde x^2z)$ vanishes except the same-index triple, so
\[
\E(\widetilde x^2z)
=\sum_jw_j^3
\E\!\left[\widetilde R_j^2(\widetilde R_j-\E\widetilde R_j)\right]
=(m_3-m_2)s_3+o(\alpha^2),
\]
where the error follows from the $r=1,2,3$ tail bounds.  Since $b_\alpha=o(\alpha^2)$ and $V_A=O(\alpha)$,
\[
\E\bigl[\widetilde x^2(z+b_\alpha)\bigr]
=\E(\widetilde x^2z)+o(\alpha^2).
\]
This proves \eqref{eq:m20}--\eqref{eq:m21}.

\paragraph{Fourth moment of $x$.}
For independent centered variables $w_jA_j$,
\[
\E\widetilde x^4
=3V_A^2+D_4,
\]
where
\[
D_4=\sum_jw_j^4
\left(\E\widetilde R_j^4-3(\E\widetilde R_j^2)^2\right).
\]
If $p<4$ (equivalently, $0<\delta<1$), \eqref{eq:trunc-master} gives
\[
\sum_jw_j^4\E\widetilde R_j^4
=O\!\left(\alpha^{p-1+\gamma(4-p)}\right).
\]
Writing $p=3+\delta$, the exponent is
\[
p-1+\gamma(4-p)=2+\delta+\gamma(1-\delta)>2
\]
because $1-\delta>0$, so this term is $o(\alpha^2)$.  More explicitly,
\[
\frac{1}{\alpha^2}\sum_jw_j^4\E\widetilde R_j^4
=O\!\left(\alpha^{\delta+\gamma(1-\delta)}\right)\to0.
\]
This estimate is not uniform as $\delta\downarrow0$: its constant contains $\E R^{3+\delta}$ and its decay exponent can become arbitrarily small.  If $p\ge4$, the finite fourth moment gives $O(s_4)=O(\alpha^3)$.  Also $\sum_jw_j^4(\E\widetilde R_j^2)^2=O(s_4)$.  Thus $D_4=o(\alpha^2)$ and, by \eqref{eq:VA-VB},
\[
\E\widetilde x^4=3m_2^2s_2^2+o(\alpha^2).
\]
This proves \eqref{eq:m40}.  In particular, no untruncated fourth moment is invoked when $3<p<4$.

\paragraph{The mixed fourth moment.}
Because $\E(A_jB_j)=0$ by the independent Rademacher sign, the pair-partition formula gives
\[
\E(\widetilde x^2z^2)=V_AV_B+D_{22},
\]
where $D_{22}$ is a same-index fourth-degree diagonal.  Its absolute value is bounded by a constant multiple of
\[
\sum_jw_j^4\left(\E\widetilde R_j^4
+\E\widetilde R_j^2+1\right)=o(\alpha^2)
\]
by the same diagonal estimate used for $D_4$.  Thus $D_{22}=o(\alpha^2)$ for each fixed law satisfying the assumptions, with the same non-uniformity near the third-moment boundary.  Hence
\[
\E(\widetilde x^2z^2)
=m_2(m_2-1)s_2^2+o(\alpha^2).
\]
The deterministic truncation bias is harmless:
\[
\E\bigl[\widetilde x^2(z+b_\alpha)^2\bigr]
=\E(\widetilde x^2z^2)
+2b_\alpha\E(\widetilde x^2z)
+b_\alpha^2V_A,
\]
and the last two terms are $o(\alpha^2)$ by \eqref{eq:bias}.  After centering,
the singleton-$y$ configurations vanish; their residual contribution is
contained in $b_\alpha=o(\alpha^2)$.  Equation \eqref{eq:m22} follows.

\paragraph{Fifth-degree remainder.}
Rosenthal's inequality applied to the centered truncated sums gives
\[
\E|\widetilde x|^5+\E|z|^5
\le C\left[O(\alpha^{5/2})
+\sum_jw_j^5\E\widetilde R_j^5+O(s_5)\right].
\]
If $p<5$ (equivalently, $0<\delta<2$), Lemma 1 bounds the diagonal term by
\[
O\!\left(\alpha^{p-1+\gamma(5-p)}\right)
=O\!\left(\alpha^{2+\delta+\gamma(2-\delta)}\right)
=o(\alpha^2),
\]
and if $p\ge5$ it is $O(s_5)=O(\alpha^4)$.  Together with $b_\alpha=o(\alpha^2)$ this yields
\[
\E(|\widetilde x|+|z+b_\alpha|)^5=o(\alpha^2).
\]
The indicator-removal argument at the start of the proof then gives \eqref{eq:m5}.
\end{proof}

\section{Proof of Theorem 1}
By Lemma 4, the complement of $G_\alpha$ contributes $o(\alpha^2)$ to $\E L_\alpha^2$.  On $G_\alpha$, use \eqref{eq:taylor}, Lemma 5, and \eqref{eq:m5}:
\[
\frac14\E L_\alpha^2
=m_2s_2-2(m_3-m_2)s_3
+3m_2(m_2-1)s_2^2
+2m_2^2s_2^2
+o(\alpha^2).
\]
Since
\[
s_2=\frac\alpha2+\frac{\alpha^2}{4}+O(\alpha^3),
\qquad
s_3=\frac{\alpha^2}{3}+O(\alpha^3),
\qquad
s_2^2=\frac{\alpha^2}{4}+O(\alpha^3),
\]
we obtain
\[
\frac14\E L_\alpha^2
=\frac{m_2}{2}\alpha
+\frac{15m_2^2+2m_2-8m_3}{12}\alpha^2
+o(\alpha^2).
\]
Symmetry gives $\E L_\alpha=0$, so \eqref{eq:var-exp} follows.  Finally,
\[
\Var(L_{1/n})=\frac{C_F^2}{n}\left(1+\frac{K_F}{n}+o(n^{-1})\right),
\]
and inversion gives \eqref{eq:g-shift}.

\section{Symbolic coefficient check}
The accompanying script
\texttt{supplement/moment\_expansion\_audit.py} reconstructs the relevant
formal moment and cumulant expressions with exact SymPy arithmetic.  In
particular, it independently reproduces
\[
\frac{15m_2^2+2m_2-8m_3}{3}
\]
as the coefficient of $\alpha^2$ in \eqref{eq:var-exp}, together with the
Laplace, Rademacher, Gaussian, and Student-law substitutions used elsewhere in
the project.  The script is run by continuous integration.

This reproduces the finite coefficient calculation but is not a second proof
of Theorem 1.  The theorem additionally depends on the
localization, logarithmic-moment, indicator-removal, and Taylor-remainder
arguments in the preceding sections.

\section{\texorpdfstring{Student-$t$}{Student-t} corollary}
\begin{corollary}[Variance-normalized Student-$t_\nu$, $\nu>3$]
Let $\varepsilon$ be a centered Student-$t_\nu$ random variable, rescaled to variance one, with $\nu>3$.  Then all assumptions of Theorem 1 hold (choose any $p\in(3,\nu)$).  Moreover
\begin{equation}\label{eq:Cnu}
\boxed{
C_\nu
=\sqrt{\frac{2\pi}{\nu-2}}\,
\frac{\Gamma(\nu/2)}{\Gamma((\nu-1)/2)} }
\end{equation}
and
\begin{equation}\label{eq:Knu}
\boxed{
K_\nu
=\frac54 C_\nu^2-\frac73-\frac{8}{3(\nu-3)}. }
\end{equation}
Hence, for each fixed $\nu>3$,
\[
\boxed{
 g_\nu^*(n)=n-K_\nu+o(1).}
\]
The remainder is pointwise in $\nu$; no uniformity as $\nu\downarrow3$ is asserted.
\end{corollary}

\begin{proof}
For raw Student-$t_\nu$,
\[
\E|T|^r
=\nu^{r/2}\frac{\Gamma((r+1)/2)\Gamma((\nu-r)/2)}{\sqrt\pi\,\Gamma(\nu/2)},
\qquad r<\nu.
\]
Scale cancels from $R=|\varepsilon|/\E|\varepsilon|$.  Formula \eqref{eq:Cnu} follows from $C_\nu^2=2m_2$.  Also
\[
\frac{m_3}{m_2}
=\frac{\E|T|^3}{\E|T|\,\E T^2}
=\frac{2(\nu-2)}{\nu-3}.
\]
Substitution into \eqref{eq:K} gives \eqref{eq:Knu}.  Student-$t$ has finite logarithmic moments near zero and at infinity, so the logarithmic assumption is automatic.  For a given $\nu>3$, one may choose any fixed $\delta$ with $0<\delta<\nu-3$.  The constants necessarily deteriorate as $\nu\downarrow3$, because the available $\delta$ shrinks and $\E|T|^{3+\delta}$ diverges as $3+\delta\uparrow\nu$.
\end{proof}

Selected values are
\begin{center}
\begin{tabular}{ccc}
\toprule
$\nu$ & $C_\nu$ & $K_\nu$\\
\midrule
$4$ & $2$ & $0$\\
$4.5$ & $1.954366$ & $0.663323$\\
$5$ & $1.923825$ & $0.959710$\\
$5.1013\ldots$ & $1.918833$ & $1$\\
$6$ & $1.885618$ & $11/9$\\
$8$ & $1.847521$ & $7/5$\\
$\infty$ (Gaussian) & $\sqrt\pi$ & $1.593657\ldots$\\
\bottomrule
\end{tabular}
\end{center}
The value $\nu=4$ is a genuine cancellation point: $K_4=0$, so the first constant correction vanishes and
\[
g_4^*(n)=n+o(1).
\]
This does not mean that the finite-$n$ normalization is exactly $n$; it means only that the constant-shift term is zero.  At the opposite edge, $K_\nu$ is not bounded as $\nu\downarrow3^+$ because of the pole $-8/[3(\nu-3)]$.  This divergence is consistent with, and warns against, any uniform interpretation of the $o(1)$ remainder across the full family $\nu>3$.

The familiar proxy $\sqrt{n-1}$ is therefore not a Student-$t$ law in general.  Within this family it is first-order variance-correct only at a root of $K_\nu=1$; numerically,
\[
\nu_*\approx5.1013.
\]
Thus the appearance of $n-1$ corresponds to a particular tail thickness, not to a universal geometric law.

\section{Scope and limitations}
This note does not address the boundary $\nu=3$ or the regime $2<\nu<3$.
The shrinking-cutoff proof requires a moment strictly above order three.

The note also does not derive higher-order coefficients under polynomial tails.
It establishes only the constant shift $K_F$, generated by the $\alpha^2$ term
in the variance.  Corrections of order $1/n$ and $1/n^2$ would require control
of progressively higher diagonal moments and their threshold cases.

\section{Conclusion}
For every fixed Student-$t_\nu$ law with $\nu>3$, the first-order
variance-normalizing length has the form
\[
g_\nu^*(n)=n-K_\nu+o(1),
\]
with $K_\nu$ given by \eqref{eq:Knu}.  This extends the constant-shift
calculation beyond exponential tails, but the remainder is pointwise in
$\nu$ and the result does not extend uniformly to the third-moment boundary.

\end{document}
